\section{Introduction}\label{sec:intro}

\paragraph{Payment Channel Networks}
Bitcoin~\cite{nakamoto2008}, the leading cryptocurrency by market
capitalisation, is costly, has low bandwidth (a few dozen transactions per
second) and high latency. Payment channels let two users pool funds, then
perform arbitrarily many payments quickly and cheaply via balance updates,
relying on the blockchain only to resolve disputes; revocation ensures that no
honest party loses money. As a channel per pair of users is often not cost
effective, payment channel networks route funds transitively across channels.
The most popular is perhaps the Lightning Network~\cite{poon2016}, with about
\$170M locked~\cite{oneml2026}. Due to the sheer complexity of its
specification~\cite{bolts}, it suffers from many attacks~\cite{tikhomirov2020},
\cite{malavolta2019}, \cite{aumayr2021}. Despite promising ongoing
efforts~\cite{grundmann2026}, \cite{litos2026}, it still lacks a
machine-checked model of both the channel layer and multi-hop routing for
unboundedly many parties, channels, updates and payments and for routes of any
length (\cite{grundmann2026} model-checks finite instances in TLA+). This paper
attempts to alleviate that issue.
% TODO: keep this sentence consistent with the characterisation of [9] in Sec. VII
% (scope and tool of [9] still to be stated there).

\paragraph{Protocol Verification}
Designing and analysing cryptographic protocols is notoriously difficult and
error-prone, as the attacks and tentative fixes on the Lightning Network
illustrate. In the computational model, inherited from cryptography, attackers
and participants are polynomial-time algorithms and results are typically
probabilistic; its tools, such as \textsc{CryptoVerif}~\cite{blanchet2006},
\textsc{EasyCrypt}~\cite{barthe2011}, \textsc{Squirrel}~\cite{baelde2021} and
\textsc{CryptoVampire}~\cite{jeanteur2024}, are all brittle, hard to use or
narrow in scope. The symbolic (\emph{Dolev-Yao}) model~\cite{dolevyao1983}
gives the attacker infinite computational power but limits it to a fixed term
algebra, abstracting away from cryptography to focus on the logic of a
protocol, where much attacker activity is found; it enjoys
successful automated verifiers, \proverif~\cite{blanchet2001} and
\tamarin~\cite{meier2013}. Mechanised proofs cover widely used protocols such
as \textsc{WireGuard}~\cite{lipp2019} and TLS~1.3~\cite{cremers2017}, the
latter in \tamarin, which proves trace properties for unboundedly many sessions
and states assumptions explicitly, as trace restrictions.
% [17] is a computational (CryptoVerif) proof of WireGuard, not a Tamarin model.

\paragraph{Contributions}
All results are checked with \tamarin~1.10 and reproduced by \texttt{make all}
(Sec.~\ref{sec:verified}); all-traces statements hold for every trace
satisfying the model's restrictions (Sec.~\ref{sec:method}).
\begin{itemize}[\setlength{\IEEElabelindent}{0pt}]
\item \emph{One theory, one adversary.} \code{lightning.spthy} covers the
  channel layer (opening, unboundedly many updates, revocation, cooperative and
  unilateral closing, punishment) and multi-hop HTLC routing, with one generic
  rule, \code{Forward_HTLC}, for every hop of routes of any length. The
  Dolev-Yao adversary can also compromise any long-term key at any time, and
  two compromised parties on a payment's route can relay its preimage around
  the parties between them (\code{AttackerSkip}; Sec.~\ref{sec:threat}). The layers share
  the channel topology and a per-channel slot; their contract, that each HTLC
  is resolved at most once, is a named restriction (\code{OneResolutionPerHTLC})
  justified informally, not a composition theorem (Sec.~\ref{sec:method}).
\item \emph{Safety and liveness.} Among 32 all-traces safety lemmas: a party is
  punished only if it published a revoked state or its key was compromised; a
  sender settles only against the receiver's invoice, unless the receiver's
  key was compromised; preimages stay secret (our model never sends them); no
  HTLC is both redeemed and refunded. Under three liveness restrictions, T1 to
  T3, an intermediary whose outgoing HTLC is redeemed has its incoming one
  refunded only if its key was compromised before that HTLC timed out, and
  payment atomicity holds without key compromise. The restrictions
  remove every violating trace, including traces in which an uncompromised
  forwarder could not have acted in time; our results say nothing about such
  traces. Ten exists-trace lemmas show that the protocol runs and that attacks
  with compromised keys are reachable (Sec.~\ref{sec:verification}). T3
  (\code{IntermediaryMustClaim}) is necessary: without it, \tamarin{} finds a
  compromise-free trace in which an intermediary's outgoing HTLC is
  redeemed while its incoming one is refunded. None of the five liveness lemmas
  we check needs T1, and neither intermediary-safety lemma needs T2, from which
  the CLTV-delta ordering of timeouts (\code{ForwardTimeGap}) is derived
  (Sec.~\ref{sec:timing}).
\item \emph{The wormhole: trace and theorem.} \tamarin{} checks (by
  replaying a saved proof) a five-party trace where the payment settles
  while an honest intermediary is refunded on its outgoing HTLC, earning no
  fee (Sec.~\ref{sec:wormhole}). Conversely, if the sender settles, a forwarder
  on its route is refunded on its outgoing HTLC only if some long-term key is
  compromised: at each of the first six hops of a route of any length under
  the liveness restrictions and, if the payment hash is offered once, at every
  hop without them (so also with them). Neither result says whose key is
  compromised, or when (Sec.~\ref{sec:bypass}).
\item \emph{Countermeasure.} On a fixed route of three intermediaries, per-hop
  locks derived with secret tweaks (a symbolic stand-in for point time-locked
  contracts, PTLCs~\cite{malavolta2019}) prevent bypassing the honest middle
  hop unless the sender is compromised; plain HTLCs admit it
  (Sec.~\ref{sec:ptlc}).
\end{itemize}
Sec.~\ref{sec:overview} recalls the protocol,
Secs.~\ref{sec:symbolic}--\ref{sec:bypass} give our model and results,
Secs.~\ref{sec:related}--\ref{sec:limitations} related work and limitations.
